Consequently, \(\mathrm{P}(g_{\mathfrak{J}})\) is in bijection with the set of
morphisms of algebras \(\mathcal{A}\to\mathcal{O}_{Y}\) which reduce according
to \(\varepsilon_{\mathfrak{J}}\), and via this bijection, \(g\) corresponds
to \(\varepsilon\).

Set \(\mathcal{M}=g^{*}(\Omega^{1}_{X/S})\). Then
\(\Hom_{\mathcal{O}_{Y}\textup{-alg.}}(\mathcal{A},\mathcal{O}_{Y})\) is
identified with the set of \(\mathcal{O}_{Y}\)-morphisms
\(\psi\colon\mathcal{M}\to\mathcal{O}_{Y}\) such that \(\Im(\psi)\) is an ideal
of square zero, and one is interested in those which induce the zero
morphism \(\mathcal{M}\to\mathcal{O}_{Y_{\mathfrak{J}}}
=\mathcal{O}_{Y}/\mathfrak{J}\), \ie which map \(\mathcal{M}\) into
\(\mathfrak{J}\). Conversely, since \(\mathfrak{J}^{2}=0\), every
\(\mathcal{O}_{Y}\)-morphism \(\phi\colon\mathcal{M}\to\mathfrak{J}\) comes
from a (unique) morphism of algebras \(\mathcal{A}\to\mathcal{O}_{Y}\),
reducing according to \(\varepsilon_{\mathfrak{J}}\). Finally, one has
\[
\Hom_{\mathcal{O}_{Y}}(g^{*}(\Omega^{1}_{X/S}),\mathfrak{J})
=\Hom_{\mathcal{O}_{Y_{\mathfrak{J}}}}(g_{\mathfrak{J}}^{*}(\Omega^{1}_{X/S}),\mathfrak{J})
\]
since \(\mathfrak{J}^{2}=0\) (\cf the proof of~(b) already seen). One
therefore obtains a bijection
\[
\mathrm{P}(g_{\mathfrak{J}})
\simeq
\Hom_{\mathcal{O}_{Y_{\mathfrak{J}}}}(g_{\mathfrak{J}}^{*}(\Omega^{1}_{X/S}),\mathfrak{J})
\]
by which \(g\) corresponds to the zero morphism.

For every \(m\in\Hom_{\mathcal{O}_{Y_{\mathfrak{J}}}}(g_{\mathfrak{J}}^{*}(\Omega^{1}_{X/S}),\mathfrak{J})\),
let us denote by \(m\cdot g\) the element of \(\mathrm{P}(g_{\mathfrak{J}})\)
associated with \(g\) and \(m\) by the preceding bijection. One has already
seen that \(0\cdot g=g\); it remains to see that
\[
(0.1.8\ (*))\qquad m'\cdot(m\cdot g)=(m+m')\cdot g.
\]
This is verified locally.%
{\renewcommand{\thefootnote}{(14)}\nde{What precedes is reproduced almost
word for word from SGA~1, III~5.1; what follows has been spelled out.}}
Indeed, the two preceding morphisms \(Y\to X\) induce the same continuous
map as \(g\) between the underlying topological spaces; it therefore
suffices to check that for every affine open \(U=\Spec(A)\) of \(X\) lying
over an affine open \(\Spec(\Lambda)\) of \(S\), and every affine open
\(V=\Spec(B)\) of \(g^{-1}(U)\), they induce the same morphism of
\(\Lambda\)-algebras \(A\to B\).

Let \(J=\Gamma(V,\mathfrak{J})\) and let \(\phi\), \(\psi\) and \(\eta\) be the
morphisms \(A\to B\) induced by \(g\), \(m\cdot g\) and \(m'\cdot(m\cdot g)\)
respectively; they coincide modulo \(J\). One can write in a unique way
\(\psi=\phi+D\) (\resp \(\eta=\psi+D'\)), where \(D\) (\resp \(D'\)) is an
element of
\[
\operatorname{Der}_{\phi}(A,J)
=\{\delta\in\Hom_{\Lambda}(A,J)\mid
\delta(ab)=\phi(a)\delta(b)+\phi(b)\delta(a)\}
\]
(\resp \(\operatorname{Der}_{\psi}(A,J)\)). But
\(\operatorname{Der}_{\phi}(A,J)=\operatorname{Der}_{\psi}(A,J)\) since
\(J^{2}=0\), and both are identified with
\[
\Hom_{B/J}(\Omega^{1}_{A/\Lambda}\otimes_{A}B/J,\,J),
\]
and via this identification \(D\) corresponds to \(m\) and \(D'\) to \(m'\).
Then \(\eta=\phi+D+D'\) and \(D+D'\) corresponds to \(m+m'\), whence the
equality~\((*)\).\qed

\begin{corollary}{0.1.9}
\label{III.0.1.9}
{\renewcommand{\thefootnote}{(15)}\nde{\normalfont One has added this
corollary of SGA~1, III~5.1, which proves point~(i) of proposition~0.2
further on.}}%
Let \(X\) be an \(S\)-scheme; let us take up again the notations of~0.1.5.
Then \(X\) is equipped with a (left) operation of the abelian
\(X^{+}\)-group \(L_{X}\), which makes \(X\) into a formally
principal homogeneous object under \(L_{X}\) over \(X^{+}\), \ie
one has an isomorphism of \(X^{+}\)-functors:
\[
L_{X}\times_{X^{+}}X
\isomto
X\times_{X^{+}}X
\]
(defined set-theoretically by \((m,x)\mapsto(x,\,m\cdot x)\)).
\end{corollary}

\begin{proof}
Let \(i_{0}\) be the immersion \(X_{0}\hookrightarrow X\). Let us first note
that, since \(X_{0}=X\times_{S}S_{0}\), one has
\(i_{0}^{*}(\Omega^{1}_{X/S})\simeq\Omega^{1}_{X_{0}/S_{0}}\)
(\cf EGA~IV,~16.4.5).

Let \(Y\) be an \(X^{+}\)-scheme, given by an \(S\)-morphism
\(g_{\mathcal{J}}\colon Y_{\mathcal{J}}\to X\), and let
\(g_{0}\colon Y_{0}\to X_{0}\) be the morphism obtained by base change.
By~0.1.8, if \(\Hom_{X^{+}}(Y,X)\) is non-empty, it is a principal
homogeneous set under the group
\[
\Hom_{\mathcal{O}_{Y_{0}}}(g_{0}^{*}i_{0}^{*}(\Omega^{1}_{X/S}),\,\mathcal{J}\mathcal{O}_{Y}),
\]
which is identified with
\(\Hom_{\mathcal{O}_{Y_{0}}}(g_{0}^{*}(\Omega^{1}_{X_{0}/S_{0}}),\,\mathcal{J}\mathcal{O}_{Y})
=L_{X}(Y)\). One therefore has a bijection
\[
L_{X}(Y)\times\Hom_{X^{+}}(Y,X)
\isomto
\Hom_{X^{+}}(Y,\,X\times_{X^{+}}X)
\]
given by \((m,g)\mapsto(g,\,m\cdot g)\). Let us show that this is ``functorial
in \(Y\)''.

Let \(f\colon Z\to Y\) be a morphism of \(S\)-schemes. The point is to show
that the diagram below is commutative:
\[
\xymatrix@C=3.2em@R=2.4em{
L_{X}(Y)\times\Hom_{X^{+}}(Y,X)
  \ar[r]^{\sim}
  \ar[d]_{L_{X}(f)\times f}
& \Hom_{X^{+}}(Y,\,X\times_{X^{+}}X)
  \ar[d]^{f\times f}
\\
L_{X}(Z)\times\Hom_{X^{+}}(Z,X)
  \ar[r]^{\sim}
& \Hom_{X^{+}}(Z,\,X\times_{X^{+}}X).
}
\]
If \(\Hom_{X^{+}}(Y,X)=\emptyset\), there is nothing to show. It therefore
suffices to see that, for every \(S\)-morphism \(g\colon Y\to X\) extending
\(g_{\mathcal{J}}\) and every \(m\in L_{X}(Y)\), one has:
\[
(0.1.9\ (*))\qquad (m\cdot g)\circ f=L_{X}(f)(m)\cdot(g\circ f).
\]
These two \(S\)-morphisms \(Z\to X\) coincide on \(Z_{\mathcal{J}}\) with
\(g_{\mathcal{J}}\circ f_{\mathcal{J}}\); in particular, they induce the
same continuous map as \(g\circ f\) between the underlying topological
spaces. Consequently, it suffices to see that, if \(z\in Z\), \(y=f(z)\),
\(x=g(y)\), and if \(A\), \(B\), \(C\) denote respectively the local rings
\(\mathcal{O}_{X,x}\), \(\mathcal{O}_{Y,y}\), \(\mathcal{O}_{Z,z}\), then
the morphisms \(A\to C\) induced by \((m\cdot g)\circ f\) and
\(L_{X}(f)(m)\cdot(g\circ f)\) coincide. Let us denote by \(s\) the
image of \(x\) in \(S\), \(\Lambda=\mathcal{O}_{S,s}\), and let \(J\) and
\(I\) be the ideals of \(\Lambda\) corresponding to \(\mathcal{J}\) and
\(\mathcal{I}\), and let \(\phi,\psi\colon A\to B\) and \(\theta\colon B\to C\)
be the morphisms of \(\Lambda\)-algebras induced by \(g\), \(m\cdot g\) and
\(f\). Then \(m\) induces an element \(D\) of
\[
\Hom_{B/IB}(\Omega^{1}_{A/\Lambda}\otimes_{A}B/IB,\,JB)
=\operatorname{Der}_{\Lambda}(A,JB)
\]
and one has \(\psi=\phi+D\); hence \((m\cdot g)\circ f\) corresponds to
\(\theta\circ\psi=\theta\circ\phi+\theta\circ D\). Now, one has seen in~0.1.5
that \(\theta\circ D\) is the image of \(L_{X}(f)(m)\) in
\[
\Hom_{C/IC}(\Omega^{1}_{A/\Lambda}\otimes_{A}C/IC,\,JC)
=\operatorname{Der}_{\Lambda}(A,JC);
\]
consequently, \(\theta\circ\phi+\theta\circ D\) is the image of
\(L_{X}(f)(m)\cdot(g\circ f)\). This proves the equality
\((0.1.9\ (*))\).
\end{proof}

\begin{corollary}{0.1.10}
\label{III.0.1.10}
{\renewcommand{\thefootnote}{(16)}\nde{\normalfont One has spelled out the
``functoriality in \(X\)'' of \(L_{X}\), in particular point~b)
below.}}%
a) \(L_{X}\) depends functorially on \(X\): for every
\(S\)-morphism \(f\colon X\to W\), there exists an \(S\)-morphism
\(L_{f}\colon L_{X}\to L_{W}\) which is a morphism
of abelian groups ``over \(f^{+}\)'', \ie the diagram
\[
\xymatrix@C=4em@R=1.8em{
L_{X} \ar[r]^{L_{f}} \ar[d]
& L_{W} \ar[d]
\\
X^{+} \ar[r]^{f^{+}}
& W^{+}
}
\]
is commutative and, for every \(Y\to X^{+}\),
\[
L_{f}(Y)\colon
\quad\Hom_{X^{+}}(Y,L_{X})\longrightarrow\Hom_{W^{+}}(Y,L_{W})
\]
(where \(Y\) is over \(W^{+}\) via \(f^{+}\)) is a morphism of abelian groups.

b) Moreover, the following diagram is commutative:
\[
\xymatrix@C=4em@R=2.2em{
L_{X}\times_{X^{+}}X
  \ar[r]^{\sim}
  \ar[d]_{L_{f}\times f}
& X\times_{X^{+}}X
  \ar[d]^{f\times f}
\\
L_{W}\times_{W^{+}}W
  \ar[r]^{\sim}
& W\times_{W^{+}}W.
}
\]
\end{corollary}

\begin{proof}
a) \(L_{f}\) is induced by the morphism of
\(\mathcal{O}_{X_{0}}\)-modules
\(f_{X_{0}/W_{0}/S_{0}}\colon
f_{0}^{*}(\Omega^{1}_{W_{0}/S_{0}})\to\Omega^{1}_{X_{0}/S_{0}}\)
(\cf~0.1.7~b)): for every \(X^{+}\)-scheme \(Y\), given by an \(S\)-morphism
\(g_{\mathcal{J}}\colon Y_{\mathcal{J}}\to X\), one has a commutative
diagram, functorial in \(Y\):
\[
\xymatrix@C=2.6em@R=2.6em{
\Hom_{\mathcal{O}_{Y_{0}}}(g_{0}^{*}(\Omega^{1}_{X_{0}/S_{0}}),\,\mathcal{J}\mathcal{O}_{Y})
  \ar[r]^{L_{f}(Y)}
  \ar[d]
& \Hom_{\mathcal{O}_{Y_{0}}}(g_{0}^{*}f_{0}^{*}(\Omega^{1}_{W_{0}/S_{0}}),\,\mathcal{J}\mathcal{O}_{Y})
  \ar[d]
\\
\{g_{\mathcal{J}}\}
  \ar[r]
& \{f_{\mathcal{J}}\circ g_{\mathcal{J}}\}
}
\]
where \(L_{f}(Y)\) is the map
\(\psi\mapsto\psi\circ g_{0}^{*}(f_{X_{0}/W_{0}/S_{0}})\), which is indeed a
morphism of abelian groups.%
{\renewcommand{\thefootnote}{(17)}\nde{One will note that \(L_{f}\)
depends only on \(f_{0}\colon X_{0}\to W_{0}\).}}

Let us prove~(b). Let \(Y\) be an \(X^{+}\)-scheme; if
\(\Hom_{X^{+}}(Y,X)=\emptyset\) there is nothing to show. Let therefore
\(g\in\Hom_{X^{+}}(Y,X)\); one must see that for every
\(m\in L_{X}(Y)\), one has:
\[
(0.1.10\ (*))\qquad f\circ(m\cdot g)=L_{f}(Y)(m)\cdot(f\circ g).
\]
Now, \(g\) being fixed, \(\Hom_{X}(Y,\,X\times_{X^{+}}X)\) is a subset of
\(\Hom_{X}(Y,\Delta^{(1)}_{X/S})\) and
\[
L_{X}(Y)
=\Hom_{\mathcal{O}_{Y_{0}}}(g_{0}^{*}(\Omega^{1}_{X_{0}/S_{0}}),\,\mathcal{J}\mathcal{O}_{Y})
=\Hom_{\mathcal{O}_{Y}}(g^{*}(\Omega^{1}_{X/S}),\,\mathcal{J}\mathcal{O}_{Y})
\]
a subset of \(L^{\square}_{X}(Y)\) (\cf~0.1.7); finally
\(L_{f}(Y)\) is the restriction to \(L_{X}(Y)\) of the map
\(L^{\square}_{f}(Y)\). Moreover, the bijection
\[
L_{X}(Y)\isomto\Hom_{X}(Y,\,X\times_{X^{+}}X),
\qquad m\mapsto(g,\,m\cdot g)
\]
is (the inverse of) the restriction to
\(L_{X}(Y)\subset L^{\square}_{X}(Y)\) of the bijection
\(\Hom_{X}(Y,\Delta^{(1)}_{X/S})\isomto
\{g\}\times L^{\square}_{X}(Y)\) considered in~0.1.7.1.
Consequently, the equality \((0.1.10\ (*))\) follows from
\((0.1.7\ (*))\); indeed, if one denotes by \(g'\) the \(X\)-morphism
\(Y\to\Delta^{(1)}_{X/S}\) defined by \((g,\,m\cdot g)\), then the element
of \(L_{W}(Y)\) corresponding to \((f\circ g,\,f\circ g')\) is
\(L^{\square}_{f}(Y)(m)=L_{f}(Y)(m)\), \ie one indeed has
\[
L_{f}(m)\cdot(f\circ g)=f\circ(m\cdot g).
\]
\end{proof}

\begin{lemma}{0.1.11}
\label{III.0.1.11}
Let \(X\), \(X'\) be two \(S\)-schemes. One has a commutative diagram:
\[
\xymatrix@C=4.5em@R=2.8em{
L_{X}\times_{S}L_{X'}
  \ar[r]^{\sim}
  \ar[d]
& L_{X\times_{S}X'}
  \ar[d]
\\
X^{+}\times_{S}X'^{+}
  \ar[r]^{\sim}
& (X\times_{S}X')^{+}.
}
\]
\end{lemma}

\begin{proof}
{\renewcommand{\thefootnote}{(18)}\nde{One has spelled out the proof.}}%
First, for every \(S\)-scheme \(Y\),
\(\Hom_{S}(Y,\,X^{+}\times_{S}X'^{+})\) equals
\(\Hom_{S}(Y,X^{+})\times\Hom_{S}(Y,X'^{+})\) and the latter is isomorphic to
\[
\Hom_{S}(Y_{\mathcal{J}},X)\times\Hom_{S}(Y_{\mathcal{J}},X')
=\Hom_{S}(Y,\,(X\times_{S}X')^{+});
\]
this proves that \(X^{+}\times_{S}X'^{+}\simeq(X\times_{S}X')^{+}\).

Next, let \(Y\) be a scheme over \(X^{+}\times_{S}X'^{+}\) via a morphism
\(h\colon Y_{\mathcal{J}}\to X\times_{S}X'\); set \(f=p\circ h\) and
\(g=q\circ h\), where one has denoted by \(p\), \(q\) the projections of
\(X\times_{S}X'\) to \(X\) and \(X'\). Since
\[
\Omega^{1}_{(X_{0}\times_{S_{0}}X'_{0})/S_{0}}
\simeq
p_{0}^{*}(\Omega^{1}_{X_{0}/S_{0}})\oplus q_{0}^{*}(\Omega^{1}_{X'_{0}/S_{0}})
\]
(\cf EGA~\(\mathrm{IV}_{4}\),~16.4.23), one obtains a natural isomorphism:
\begin{align*}
&\Hom_{\mathcal{O}_{Y_{0}}}(f_{0}^{*}(\Omega^{1}_{X_{0}/S_{0}}),\,\mathcal{J}\mathcal{O}_{Y})
\times
\Hom_{\mathcal{O}_{Y_{0}}}(g_{0}^{*}(\Omega^{1}_{X'_{0}/S_{0}}),\,\mathcal{J}\mathcal{O}_{Y})\\
&\qquad\simeq
\Hom_{\mathcal{O}_{Y_{0}}}(h_{0}^{*}(\Omega^{1}_{(X_{0}\times_{S_{0}}X'_{0})/S_{0}}),\,\mathcal{J}\mathcal{O}_{Y}),
\end{align*}
\ie \(L_{X}(Y)\times L_{X'}(Y)\simeq L_{X\times_{S}X'}(Y)\).
\end{proof}

\begin{remark}{0.1.12}
\label{III.0.1.12}
{\renewcommand{\thefootnote}{(19)}\nde{One has added this remark, as well as
the corollary which follows.}}%
Let \(\mathcal{C}\) be a category stable under fibered products, \(S\) an
object of \(\mathcal{C}\), \(T_{1}\), \(T_{2}\) two objects over \(S\) and,
for \(i=1,2\), \(L_{i}\) and \(X_{i}\) two objects over \(T_{i}\):
\[
\xymatrix@C=1.5em@R=1.6em{
L_{1}\ar[dr]
& & X_{1}\ar[dl]
& & L_{2}\ar[dr]
& & X_{2}\ar[dl]
\\
& T_{1}\ar[drr]
& & & & T_{2}\ar[dll]
\\
& & & S
}
\]
Then one has a natural isomorphism:
\[
(L_{1}\times_{T_{1}}X_{1})\times_{S}(L_{2}\times_{T_{2}}X_{2})
\simeq
(L_{1}\times_{S}L_{2})\times_{T_{1}\times_{S}T_{2}}(X_{1}\times_{S}X_{2}).
\]
Consequently, one deduces from the preceding lemma the:
\end{remark}

\begin{corollary}{0.1.13}
\label{III.0.1.13}
Let \(X_{1}\), \(X_{2}\) be two \(S\)-schemes. One has a commutative diagram
of isomorphisms:
\[
\xymatrix@C=1.2em@R=2.6em{
*+{L_{X_{1}\times_{S}X_{2}}\times_{(X_{1}\times_{S}X_{2})^{+}}(X_{1}\times_{S}X_{2})}
  \ar[d]_{(0.1.11)\ \simeq}
  \ar[dr]^{\simeq}
&
\\
*+{(L_{X_{1}}\times_{S}L_{X_{2}})\times_{(X_{1}^{+}\times_{S}X_{2}^{+})}(X_{1}\times_{S}X_{2})}
  \ar[r]^{\sim}_{(0.1.12)}
& *+{(L_{X_{1}}\times_{X_{1}^{+}}X_{1})\times_{S}(L_{X_{2}}\times_{X_{2}^{+}}X_{2}).}
}
\]
\end{corollary}

We can now state:

\begin{proposition}{0.2}
\label{III.0.2}
For every \(S\)-scheme \(X\), one can define a (left) operation of the
abelian \(X^{+}\)-group \(L_{X}\) on the \(X^{+}\)-object \(X\),
such that:
\begin{enumeratei}
\item this operation makes \(X\) into a formally principal homogeneous
object under \(L_{X}\) over \(X^{+}\), \ie the morphism
\[
L_{X}\times_{X^{+}}X\longrightarrow X\times_{X^{+}}X
\]
is an isomorphism of \(X^{+}\)-functors;
\item this operation is functorial in the \(S\)-scheme \(X\), \ie for every
\(S\)-morphism \(f\colon X\to W\), the following diagram is commutative:
\[
\xymatrix@C=4em@R=2.2em{
L_{X}\times_{X^{+}}X
  \ar[r]
  \ar[d]_{L_{f}\times f}
& X
  \ar[d]^{f}
\\
L_{W}\times_{W^{+}}W
  \ar[r]
& W\,;
}
\]
\item this operation ``commutes with fibered product'', \ie for all
\(S\)-schemes \(X_{1}\) and \(X_{2}\), the following diagram is commutative:
\[
\xymatrix@C=1em@R=2.8em{
*+{L_{X_{1}\times_{S}X_{2}}\times_{(X_{1}\times_{S}X_{2})^{+}}(X_{1}\times_{S}X_{2})}
  \ar[rr]
  \ar[d]_{\simeq}
  \ar[drr]^{\simeq}
& & X_{1}\times_{S}X_{2}
\\
*+{(L_{X_{1}}\times_{S}L_{X_{2}})\times_{(X_{1}\times_{S}X_{2})^{+}}(X_{1}\times_{S}X_{2})}
  \ar[rr]^{\sim}
& & *+{(L_{X_{1}}\times_{X_{1}^{+}}X_{1})\times_{S}(L_{X_{2}}\times_{X_{2}^{+}}X_{2}).}
  \ar[u]
}
\]
\end{enumeratei}
\end{proposition}

\begin{proof}
{\renewcommand{\thefootnote}{(20)}\nde{One has modified and spelled out the
original, taking into account the additions made previously; these
incorporate, spelling it out, the content of page~89 of the original.}}%
(i)~and~(ii) follow respectively from corollaries~0.1.9 and~0.1.10. To
prove~(iii), set \(\mathrm{P}(X)=L_{X}\times_{X^{+}}X\),
for every \(S\)-scheme \(X\). Then, by~(ii) applied to the projections
\(p_{i}\colon X_{1}\times_{S}X_{2}\to X_{i}\), one obtains commutative
squares
\[
\xymatrix@C=4em@R=2.2em{
\mathrm{P}(X_{1}\times_{S}X_{2})
  \ar[r]
  \ar[d]_{L_{p_{i}}\times p_{i}}
& X_{1}\times_{S}X_{2}
  \ar[d]^{p_{i}}
\\
\mathrm{P}(X_{i})
  \ar[r]
& X_{i}
}
\]
for \(i=1,2\), and hence a commutative square:
\[
\xymatrix@C=4em@R=2.4em{
\mathrm{P}(X_{1}\times_{S}X_{2})
  \ar[r]
  \ar[d]
& X_{1}\times_{S}X_{2}
\\
\mathrm{P}(X_{1})\times_{S}\mathrm{P}(X_{2})
  \ar[r]
& X_{1}\times_{S}X_{2}.
}
\]
Combining this with corollary~0.1.13, one obtains that the vertical arrow
is an isomorphism, and that one has the commutative diagram announced
in~(iii).
\end{proof}

\begin{remark}{0.3}
\label{III.0.3}
Suppose the \(X^{+}\)-scheme \(Y\) \emph{flat} over \(S\) (\cf SGA~1,~IV).
One can then write
\[
\Hom_{X^{+}}(Y,L_{X})
=\Hom_{\mathcal{O}_{Y_{0}}}(g_{0}^{*}(\Omega^{1}_{X_{0}/S_{0}}),\,
\mathcal{J}\otimes_{\mathcal{O}_{S_{0}}}\mathcal{O}_{Y_{0}}).
\]
\end{remark}

\begin{remark}{0.4}
\label{III.0.4}
Let us denote by \(\pi_{0}\colon X_{0}\to S_{0}\) the structural morphism
and suppose that there exists an \(\mathcal{O}_{S_{0}}\)-module
\(\omega^{1}_{X_{0}/S_{0}}\) such that
\(\Omega^{1}_{X_{0}/S_{0}}\simeq\pi_{0}^{*}(\omega^{1}_{X_{0}/S_{0}})\)
(the case will present itself in particular when \(X_{0}\) is an
\(S_{0}\)-group, \cf II,~4.11). If one defines a functor \(L'_{X}\)
over \(S\) by the formula
\[
(0.4.1)\qquad
\Hom_{S}(Y,L'_{X})
=\Hom_{\mathcal{O}_{Y_{0}}}(\omega^{1}_{X_{0}/S_{0}}\otimes_{\mathcal{O}_{S_{0}}}\mathcal{O}_{Y_{0}},\,
\mathcal{J}\mathcal{O}_{Y}),
\]
one then has \(\Hom_{X^{+}}(Y,L_{X})=\Hom_{S}(Y,L'_{X})\)
for every \(X^{+}\)-scheme \(Y\), \ie
\[
L_{X}=L'_{X}\times_{S}X^{+}.
\]
{\renewcommand{\thefootnote}{(21)}\nde{One has spelled out what follows.}}%
Then, since \(L_{X}\times_{X^{+}}X=L'_{X}\times_{S}X\),
the operation of \(L_{X}\) on \(X\) induces an operation of
\(L'_{X}\) on \(X\), and this operation respects the morphism
\(p_{X}\colon X\to X^{+}\); indeed, if \(Y\) is an \(S\)-scheme,
\(h\colon Y\to X\) an \(S\)-morphism and \(m\) an element of
\(L'_{X}(Y)\), then \(h\) and \(m\cdot h\) have the same restriction
to \(Y_{\mathcal{J}}\), \ie \(p_{X}(m\cdot h)=p_{X}(h)\).
\end{remark}

\begin{remark}{0.5}
\label{III.0.5}
Let us keep the hypotheses and notations of~0.4 and suppose moreover that
\(Y\) is an \(S\)-scheme \emph{flat} over \(S\). One then has
\[
\Hom_{X^{+}}(Y,L_{X})
=\Hom_{S}(Y,L'_{X})
=\Hom_{S_{0}}(Y_{0},L_{0X}),
\]
where the \(S_{0}\)-functor in abelian groups \(L_{0X}\) is defined
by the following identity (with respect to the variable \(S_{0}\)-scheme
\(T\)):
\[
(0.5.1)\qquad
\Hom_{S_{0}}(T,L_{0X})
=\Hom_{\mathcal{O}_{T}}(\omega^{1}_{X_{0}/S_{0}}\otimes_{\mathcal{O}_{S_{0}}}\mathcal{O}_{T},\,
\mathcal{J}\otimes_{\mathcal{O}_{S_{0}}}\mathcal{O}_{T}).
\]
In the notations of~II,~1, one has therefore shown that the functors
\(L'_{X}\) and \(\prod_{S_{0}/S}L_{0X}\) have the same
restriction to the full subcategory of \((\mathbf{Sch})_{/S}\) whose objects
are the \(S\)-schemes \(Y\) \emph{flat} over \(S\).
\end{remark}

\begin{remark}{0.6}
\label{III.0.6}
Let us keep the hypotheses and notations of~0.5%
{\renewcommand{\thefootnote}{(22)}\nde{One has added the hypothesis which
follows, implicit in the original.}}
and suppose moreover that there exists a section \(\varepsilon_{0}\) of
\(\pi_{0}\colon X_{0}\to S_{0}\); one then has
\(\omega^{1}_{X_{0}/S_{0}}\simeq\varepsilon_{0}^{*}(\Omega^{1}_{X_{0}/S_{0}})\).
First, one has (independently of the preceding hypothesis):
\[
\Hom_{S_{0}}(T,L_{0X})
=\Gamma\bigl(T,\,
\mathcal{H}\mathrm{om}_{\mathcal{O}_{T}}(\omega^{1}_{X_{0}/S_{0}}\otimes_{\mathcal{O}_{S_{0}}}\mathcal{O}_{T},\,
\mathcal{J}\otimes_{\mathcal{O}_{S_{0}}}\mathcal{O}_{T})\bigr).
\]
\end{remark}
